\documentclass[11pt,a4paper]{article}
\usepackage[T1]{fontenc}
\usepackage{lmodern,amsmath,amssymb,amsthm,mathtools,microtype}
\usepackage[margin=27mm]{geometry}
\usepackage[colorlinks=true,linkcolor=blue,citecolor=blue,urlcolor=blue]{hyperref}
\setlength{\emergencystretch}{2em}
\newtheorem{theorem}{Theorem}[section]
\newtheorem{lemma}[theorem]{Lemma}
\newcommand{\Q}{\mathbb Q}\newcommand{\Z}{\mathbb Z}
\DeclareMathOperator{\Tr}{Tr}\DeclareMathOperator{\Nm}{N}
\title{Intrinsic congruent-curve poles from a fixed scalar rescaling}
\author{}\date{17 November 2026}
\begin{document}\maketitle
\begin{abstract}
A change of two fixed scalar coefficients puts the normalized general
elliptic source in the prepared range of an intrinsic congruent-curve
pole equation. The physical detector weights and all point roles remain
unchanged. An exhaustive certificate at the consumed prime verifies the
only new finite-residue condition. Ordered completeness uses Burungale's
simultaneous-twist theorem, including density for all-character guard
avoidance. Under the remaining stated source premises the resulting
family defines the integers with twenty-five total variables. The source
also retains the hypotheses of the native square-cover criterion, whose
smaller realization problem is not resolved here.
\end{abstract}

\section{Scalar specification}

Write either retained general scalar in the form
\[
 \rho_L(X)=\frac{A_Lr(X)}{(1+C_Lr(X))^2},\qquad \ell=76249,
 \qquad L\ge3\text{ odd}.
\]
For the original scalar of \cite{General},
\[
 \beta=216X^4+450459,\qquad A_L=102^2\ell^L,
 \qquad C_L=17(\ell^{(L+1)/2}+1),
\]
\[
 r=\beta(\beta-1)(\beta-2401)(32\beta+343)(2\beta-1).
\]
For the aligned scalar of \cite{Aligned}, put
\[
 n=116576341082536192,\quad D_0=5951086353,\quad
 A=1296233=17\ell,\quad Z_0=12556178240,
\]
\[
 \beta=2(AX^2+Z_0)^2-D_0^2,
 \qquad r=\beta(\beta-n)(\beta-2401n)(32\beta+343n),
\]
\[
 M=r(0),\qquad A_L=102^2\ell^LM,
 \qquad C_L=102^2M(\ell^{(L+1)/2}+1).
\]
These are the exact retained labels and contacts. Define instead
\begin{equation}\label{eq:scalar}
                 \rho_L^*(X)=\frac{7A_Lr(X)}{(1+7C_Lr(X))^2}.
\end{equation}
Use \(\rho_L^*\) in every scalar and comparison-scale occurrence.
The fixed exceptional sets of the two sources are unchanged.

\begin{lemma}\label{lem:scalar}
The scalar \eqref{eq:scalar} preserves the integrality-reflecting
label, the scalar contact-height proof, and every consumed-prime
valuation. For integer inputs it has square reduced denominator and
\[
 \rho_L^*\in7\Q_2^{\times2},\qquad
 v_2(\rho_L^*)=4\text{ or }12,\qquad v_7(\rho_L^*)\ge1.
\]
For every real input, \(0<\rho_L^*<\rho_L<1/56\).
Its height estimate has the same exponents and an absolute change
of constant. For the aligned source it remains a norm from \(\Q(\sqrt2)\).
\end{lemma}
\begin{proof}
The exact ratio is
\[
 \frac{\rho_L^*}{\rho_L}
       =7\left(\frac{1+C_Lr}{1+7C_Lr}\right)^2.
\]
Outside the fixed coefficient supports and seven, this preserves the
odd contact primes. At a positive odd contact of \(r\) the two
denominators are units, so it preserves the required odd valuation
itself. The fixed coefficient supports have exactly the treatment
already specified in the source's height proof. The scalar-height
bound uses the same polynomial degrees; multiplying the fixed
coefficients by seven changes only its absolute constant.

At \(\ell\), the scalar coefficient \(C_L\) is a unit. If the
input has negative \(\ell\)-adic valuation, the source proves
\(v_\ell(r)<0\), and both old and new denominators have valuation
\(v_\ell(r)\). If the input is \(\ell\)-integral, the exact
residue certificate accompanying this paper verifies
\[
            r(X)\not\equiv0,\qquad 1+7C_Lr(X)\not\equiv0\pmod\ell
\]
for all \(76249\) residues \(X\). Here \(C_L\equiv17\) in the
original case and \(C_L\equiv102^2M\) in the aligned case,
independently of \(L\). The aligned contact value is itself
constant modulo \(\ell\), because \(A=17\ell\).
Thus the complete consumed-prime valuation is unchanged, not just
its parity; in particular it is positive odd and at least \(L\).

For integer inputs, put \(D=1+7C_Lr\). One has \(\gcd(r,D)=1\).
In the original scalar the remaining numerator primes are
\(2,3,7,17,\ell\). The denominator is a unit at two because
\(C_L\) is even, at three because \(\beta\equiv0\pmod3\) and
therefore \(r\equiv0\pmod3\), at seventeen because \(17\mid C_L\),
and at seven by its displayed form. At \(\ell\) this is the
residue certificate just proved. For the aligned scalar, every
prime dividing \(102M\) divides \(C_L\), and again seven and
\(\ell\) give units by the preceding arguments. Thus in both cases
\(\gcd(7A_Lr,D)=1\). The reduced denominator in
\eqref{eq:scalar} is consequently the square \(D^2\), with no
odd cancellation. The numerator is divisible by seven.
At two the old and new denominators are units by the original
local preparation, and multiplying \(C_L\) by an odd integer does
not alter this fact. The displayed square ratio gives the claimed
two-adic class and unchanged valuation.

On the real line both source specifications have \(C_Lr\ge1\).
For \(t\ge1\),
\[
 \frac{7(1+t)^2}{(1+7t)^2}\le\frac7{16}<1.
\]
The original scalar bound is \(\rho_L<36/\ell<1/56\), and
the aligned bound is smaller. Finally
\(7=\Nm(3+\sqrt2)\), so the square ratio also preserves the
aligned norm assertion.
\end{proof}

\section{Detectors and pole equations}

Use either the normalized general tests of \cite{Reweighted} or the
order-two revision of \cite{Precision}. Keep the common slope, all
nine point roles, the two contact functions \(R_j\), and the
unsquared scales \(a_0=1,a_1=1+\rho_L^*T^2\).
The physical weights are the established ones:
\[
 u_j=\rho_L^*R_j,\qquad w_j=\frac{u_j}{(1+u_j)^2},
 \qquad b_j=\frac{4w_j}{\Tr(W_j^{-1})^2}.
\]
The detector \(W_j\), with its chosen field and scalar orders,
uses this weight in every denominator. Keep all point, difference,
scalar, contact and trace guards.

For each block take two rational coordinates \(v_j,s_j\) and
the finite equation
\begin{equation}\label{eq:pole}
 Y_j^2=P_{b_j}(v_j,s_j),\qquad
 P_d(v,s)=d(d+v^2)^2s^4-d(d-v^2)^2(1-d)^2.
\end{equation}
The intrinsic pole theorem of \cite{Intrinsic} proves, without
guards, that this equation excludes negative odd valuations of
\(b_j\) at every prime. Its proof uses the two distinct odd
valuations of the radicand terms at a hypothetical negative odd
valuation of \(b_j\).
Hence the source's local extraction and height soundness are
unchanged by Lemma~\ref{lem:scalar}.

\section{Ordered completeness and all-character guards}

Fix an integer input and a fixed odd \(L\ge3\). After fixing
\(\rho_L^*\), form the temporary finite preparation and avoidance
set including its support and the support of \(1+\rho_L^*\).
This is not an enlargement of the global exceptional set used
in the height proof. Choose the reference, scalar and root points
in the retained order, and fix the comparison scales before
applying the prime-pattern theorem.

All fixed local scalar functions can be made as small as needed.
At a selected contact the old normalizer and reciprocal-trace
proofs still give positive odd valuation of \(b_j\). At fixed
preparation primes the contact functions tend to one and the
detectors to one, preserving the scalar's prescribed square
classes. The unchanged even-degree contact factorization shows
that every other valuation is even. Thus both \(b_j\) have square
reduced denominator.

At two the preparation gives
\[
          b_j\in7\Q_2^{\times2},\qquad
          v_2(b_j)=v_2(\rho_L^*)\ge4,
\]
since the reciprocal trace is sufficiently close to two.
At seven \(u_j\) has positive valuation and \(1+u_j\) is a unit;
the reciprocal trace is close to two and is a unit. In particular
no negative odd valuation is introduced there.

At the selected real character choose \(R_j<2\) and both real
detector slots in \([1,2)\). The reciprocal trace exceeds one,
and \(0<w_j<u_j<2\rho_L^*\). Therefore
\[
                    0<b_j<8\rho_L^*<1/7<1.
\]
These are precisely the prepared-input hypotheses of the intrinsic
theorem, which follows from Burungale's Theorem B. It supplies a
Zariski-dense set of rational witness pairs \((v_j,s_j)\) for each
actual input.

At every geometric point-algebra character all detector factors
are positive and the traces nonzero, so \(b_j\) is a unit. Other
characters need not satisfy \(b_j<1\), and may take the value one.
No such restriction is needed in \eqref{eq:pole}.
For each character value \(c\ne0\), the polynomial \(P_c(v,s)\)
is nonzero: its leading coefficient in \(s\) is \(c(c+v^2)^2\).
Its point-algebra norm is consequently a nonzero polynomial.
The intrinsic density allows the actual witness pairs to avoid
this norm and every retained denominator norm simultaneously.
The cleared terminal radicands are thus units at every character,
and the existing one-coordinate quadratic encoder applies.

\begin{theorem}
Retain the source's group, prime-pattern, normalized contact-height
and aggregate trace-height premises, and use Burungale's Theorem B
in the intrinsic form stated in \cite{Intrinsic}.
With either scalar and either indicated set of scalar orders,
the modified source defines \(\Z\) over \(\Q\) for every sufficiently
large fixed odd \(L\). It is complete for every integer input and
every fixed odd \(L\ge3\). Its finite polynomial specification has
twenty-five total variables.
\end{theorem}
\begin{proof}
Soundness is the retained source argument with the scalar properties
of Lemma~\ref{lem:scalar} and the intrinsic all-prime pole implication.
The ordered preparation and guard argument prove completeness.
The coordinates are the input, eighteen point coordinates, one slope,
four pole coordinates and one encoder coordinate. No other free
coordinate is introduced.
\end{proof}

Every odd valuation of a physical weight outside the fixed
exceptional set is still positive. Thus the native square-cover
criterion remains applicable to this source as well. No rational
realization of that smaller cover, numerical sufficient \(L\), or
degree bound for the complete determinant is asserted here.

\begin{thebibliography}{9}\raggedright
\bibitem{General}\emph{A defining polynomial in twenty-five variables},
\nolinkurl{papers/few_variables25/few_variables25_merged_20261110_v19.tex}.
\bibitem{Aligned}\emph{A scalar norm from the detector field},
\nolinkurl{papers/pole_parameters/detector_field_norm_scalar_20261116_v52.tex}.
\bibitem{Reweighted}\emph{Quadratic reweightings of native elliptic detectors},
\nolinkurl{papers/pole_parameters/reweighted_elliptic_detectors_20261116_v53.tex}.
\bibitem{Precision}\emph{Scalar precision two with an eighth-order elliptic comparison},
\nolinkurl{papers/pole_parameters/precision_transfer_20261117_v56.tex}.
\bibitem{Intrinsic}\emph{A two-coordinate pole test of coefficient height four},
\nolinkurl{papers/pole_parameters/intrinsic_pole_20261117_v56.tex}.
\bibitem{Burungale}A.~Burungale,
\emph{A Diophantine definition of $\Z$ in $\Q$},
arXiv:2611.16130v1, 12 November 2026, Theorem B.
\url{https://arxiv.org/abs/2611.16130}.
\end{thebibliography}
\end{document}
